\chapter{Cloud Processing Options}
\label{ch:cloud}

\begin{center}
\textit{Learning Objectives: Understand cloud GPU rental options, estimate costs per scene, and choose the right tier for your project budget.}\\[0.5em]
\textbf{Prerequisites:} Module~2 (Processing Pipeline) \hfill \textbf{Time:} 30\,minutes
\end{center}

\vspace{1em}

The gaussian-toolkit pipeline is fully containerised with Docker, meaning it can run on \emph{any} cloud GPU provider without modification. This chapter presents pricing options, capacity planning, and recommendations for different project scales.

\section{Why Cloud Processing?}

Not every institution has access to dedicated GPU hardware. The pipeline requires:

\begin{itemize}
  \item \textbf{Minimum:} One NVIDIA GPU with 12\,GB VRAM (RTX 3060 or equivalent)
  \item \textbf{Recommended:} One GPU with 24--48\,GB VRAM (RTX 4090, A100, or RTX 6000 Ada)
  \item \textbf{Optimal:} Two GPUs --- one for training, one for MILo mesh extraction
\end{itemize}

Cloud rental eliminates capital expenditure and provides access to high-end hardware on demand. A complete scene capture (COLMAP + training + mesh extraction) takes approximately 30 minutes on an RTX 4090, costing as little as \pounds0.18 per scene.

\section{Cloud GPU Providers}
\label{sec:cloud-providers}

\begin{table}[htbp]
\centering
\caption{Cloud GPU provider comparison (prices as of March 2026, approximate GBP equivalents).}
\label{tab:cloud-pricing}
\small
\begin{tabular}{@{}llrrp{4.5cm}@{}}
\toprule
\textbf{Provider} & \textbf{GPU} & \textbf{\pounds/hr} & \textbf{\pounds/scene} & \textbf{Notes} \\
\midrule
RunPod        & RTX 4090    & 0.35 & 0.18 & Docker native, spot pricing available \\
Vast.ai       & RTX 4090    & 0.24 & 0.12 & Cheapest spot, community GPUs \\
Salad Cloud   & RTX 4090    & 0.20 & 0.10 & Batch processing, ComfyUI integration \\
Lambda Cloud  & A100 (40G)  & 0.88 & 0.44 & Enterprise SLA, persistent storage \\
Lambda Cloud  & H100        & 1.99 & 1.00 & Maximum performance \\
Google Colab Pro+ & A100   & 40/mo & 0.13$^\dagger$ & Unlimited A100 sessions \\
AWS g5.xlarge & A10G (24G)  & 0.81 & 0.41 & Enterprise, reserved pricing available \\
\bottomrule
\multicolumn{5}{l}{\footnotesize $^\dagger$ Assumes 300 scenes/month at 40\,\pounds/month subscription.}
\end{tabular}
\end{table}

\section{Budget Tiers}
\label{sec:budget-tiers}

\begin{recommendation}{Free Tier (\pounds0/month)}
Use Google Colab free tier (T4 GPU, 15\,GB VRAM). Sufficient for COLMAP and basic 3DGS training at reduced resolution. Processing time approximately 2$\times$ longer than RTX 4090. MILo mesh extraction not available (insufficient VRAM). Suitable for evaluation and small-scale testing.
\end{recommendation}

\begin{recommendation}{Starter Tier (\pounds50/month)}
Google Colab Pro (\pounds8/month) plus \pounds42 in RunPod credits. Covers approximately 120 scene captures per month with RTX 4090 quality. Recommended for pilot projects and proof-of-concept work. This is the tier proposed for the UoS AI KE CoP pilot.
\end{recommendation}

\begin{recommendation}{Professional Tier (\pounds200/month)}
Dedicated RunPod or Lambda Cloud allocation. Covers approximately 500+ scene captures per month. Enables MILo mesh extraction, per-object segmentation, and batch processing. Suitable for institutional deployment across multiple venues.
\end{recommendation}

\begin{recommendation}{Enterprise Tier (\pounds500+/month)}
Reserved cloud instances or on-premise GPU (RTX 4090 $\approx$ \pounds1,600 capital). At 500+ scenes/month, on-premise hardware pays for itself within 3--4 months. Recommended for sustained production workflows.
\end{recommendation}

\section{Cost Per Scene Breakdown}
\label{sec:cost-breakdown}

\begin{table}[htbp]
\centering
\caption{Detailed cost breakdown per scene on RunPod RTX 4090 (\pounds0.35/hr).}
\label{tab:cost-per-scene}
\begin{tabular}{@{}lrrr@{}}
\toprule
\textbf{Pipeline Stage} & \textbf{Time} & \textbf{GPU Cost} & \textbf{API Cost} \\
\midrule
Frame extraction       & 30\,s  & \pounds0.003 & --- \\
COLMAP reconstruction  & 5\,min & \pounds0.029 & --- \\
3DGS training (30K)    & 5\,min & \pounds0.029 & --- \\
MILo mesh extraction   & 17\,min & \pounds0.099 & --- \\
Blender texture bake   & 1\,min & \pounds0.006 & --- \\
SAM3 segmentation      & 2\,min & \pounds0.012 & --- \\
Claude Code orchestration & ---  & --- & \pounds0.05--0.20$^\ast$ \\
\midrule
\textbf{Total}         & \textbf{30\,min} & \textbf{\pounds0.18} & \textbf{\pounds0.05--0.20} \\
\bottomrule
\multicolumn{4}{l}{\footnotesize $^\ast$ Anthropic API cost depends on orchestration complexity; free with OAuth session.}
\end{tabular}
\end{table}

\section{Deployment Options}

\subsection{Option A: Fully Cloud-Based}

Upload video via the web interface running on a cloud GPU instance. The Docker container includes all dependencies. Steps:

\begin{enumerate}
  \item Rent a GPU instance (RunPod, Vast.ai, or Lambda)
  \item Deploy the gaussian-toolkit Docker image
  \item Access the web UI at \texttt{http://\textless instance-ip\textgreater:7860}
  \item Upload video, monitor pipeline, download results
\end{enumerate}

\textbf{Advantage:} Zero local hardware required. Works from any device with a browser.\\
\textbf{Disadvantage:} Video upload time depends on internet speed. 1\,GB video at 10\,Mbps = 14 minutes.

\subsection{Option B: Hybrid (Local Capture, Cloud Processing)}

Capture video locally, upload to cloud storage (S3, GCS, or direct upload), process on cloud GPU, download results.

\textbf{Advantage:} Captures at full original quality without re-encoding.\\
\textbf{Disadvantage:} Requires cloud storage management.

\subsection{Option C: On-Premise with Cloud Burst}

Run the pipeline on local hardware (university GPU workstation) for routine processing, burst to cloud for high-demand periods (e.g.\ during capture campaigns at heritage events).

\textbf{Advantage:} Lowest ongoing cost, data stays on-premise.\\
\textbf{Disadvantage:} Requires initial capital expenditure for GPU hardware.

\section{Recommendations for the UoS Pilot}

\begin{keyfinding}
For the \pounds10K AI KE CoP pilot, we recommend the \textbf{Starter Tier} (\pounds50/month cloud budget) combined with access to the University's existing SAMCT DevLab GPU (24\,GB, as specified in the bid). This provides:
\begin{itemize}
  \item On-premise processing for routine captures (DevLab GPU)
  \item Cloud burst capacity for batch processing during Lightwaves/WITW events
  \item Approximately 120 cloud-processed scenes per month
  \item Total cloud cost for 3-month pilot: \pounds150
\end{itemize}
The remaining budget supports DreamLab AI consultation time (\pounds3,600), technical staff (\pounds5,250), and equipment sourcing.
\end{keyfinding}
